\documentclass[12pt, a4paper]{ctexart}
\usepackage[margin=2cm]{geometry}
\usepackage{amsmath, amssymb}
\usepackage{enumitem}
\usepackage{fancyhdr}
\usepackage{lastpage}
\usepackage{graphicx}
\usepackage{tikz}

\setCJKmainfont{Hiragino Sans GB}
\setCJKsansfont{Heiti SC}

\pagestyle{fancy}
\fancyhf{}
\fancyhead[L]{九年级下册数学测试卷}
\fancyhead[R]{第\thepage\ 页，共 \pageref{LastPage} 页}
\renewcommand{\headrulewidth}{0.4pt}
\fancyfoot[C]{}

\setlength{\parindent}{0pt}

\begin{document}

\begin{center}
{\LARGE\bfseries 九年级下册数学测试卷}\\[0.5em]
{\large 满分：120分\quad 考试时间：120分钟}
\end{center}

\vspace{0.5em}
\noindent\textbf{注意事项：}
\begin{enumerate}[label=\arabic*., leftmargin=2em]
\item 本试卷共三大题，24小题。
\item 请将答案写在答题卡上，在本试卷上作答无效。
\item 考试结束后，请将本试卷和答题卡一并交回。
\end{enumerate}

\vspace{0.5em}

% ============================================================
% 一、选择题（共10题，每题3分，共30分）
% ============================================================
\noindent\textbf{一、选择题（本大题共10小题，每小题3分，共30分。在每小题给出的四个选项中，只有一项是符合题目要求的。）}

\begin{enumerate}[label=\arabic*.]

\item 下列函数中，$y$ 是 $x$ 的反比例函数的是（\quad）
\begin{enumerate}[label=\Alph*.]
\item $y = \dfrac{1}{x^2}$
\item $y = \dfrac{2}{x}$
\item $y = \dfrac{x}{3}$
\item $y = \dfrac{1}{x + 1}$
\end{enumerate}

\item 反比例函数 $y = \dfrac{6}{x}$ 的图象在（\quad）
\begin{enumerate}[label=\Alph*.]
\item 第一、二象限
\item 第一、三象限
\item 第二、三象限
\item 第二、四象限
\end{enumerate}

\item 已知 $\triangle ABC \sim \triangle DEF$，$\triangle ABC$ 的面积为 $4$，$\triangle DEF$ 的面积为 $9$，若 $AB = 6$，则 $DE$ 的长为（\quad）
\begin{enumerate}[label=\Alph*.]
\item $4$
\item $8$
\item $9$
\item $10$
\end{enumerate}

\item 在 $\text{Rt}\triangle ABC$ 中，$\angle C = 90^\circ$，$AC = 3$，$BC = 4$，则 $\sin A$ 的值为（\quad）
\begin{enumerate}[label=\Alph*.]
\item $\dfrac{3}{5}$
\item $\dfrac{4}{5}$
\item $\dfrac{3}{4}$
\item $\dfrac{4}{3}$
\end{enumerate}

\item 如图，小明身高 $1.6\,\text{m}$，在某时刻测得他的影长为 $2\,\text{m}$。同一时刻，一棵大树的影长为 $10\,\text{m}$，则大树的高度为（\quad）
\begin{enumerate}[label=\Alph*.]
\item $6\,\text{m}$
\item $7\,\text{m}$
\item $8\,\text{m}$
\item $10\,\text{m}$
\end{enumerate}

\item 反比例函数 $y = \dfrac{k}{x}$（$k \neq 0$）的图象经过点 $(2, -3)$，则 $k$ 的值为（\quad）
\begin{enumerate}[label=\Alph*.]
\item $6$
\item $-6$
\item $\dfrac{3}{2}$
\item $-\dfrac{3}{2}$
\end{enumerate}

\item 如图，$\triangle ABC \sim \triangle ADE$，$AD = 2$，$DB = 4$，$BC = 9$，则 $DE$ 的长为（\quad）
\begin{enumerate}[label=\Alph*.]
\item $3$
\item $4$
\item $5$
\item $6$
\end{enumerate}

\item 在 $\text{Rt}\triangle ABC$ 中，$\angle C = 90^\circ$，$\angle A = 30^\circ$，$BC = 2$，则 $AB$ 的长为（\quad）
\begin{enumerate}[label=\Alph*.]
\item $2$
\item $2\sqrt{3}$
\item $4$
\item $4\sqrt{3}$
\end{enumerate}

\item 下列四个几何体中，主视图、左视图和俯视图都是圆的是（\quad）
\begin{enumerate}[label=\Alph*.]
\item 圆柱
\item 圆锥
\item 球
\item 三棱柱
\end{enumerate}

\item 已知反比例函数 $y = \dfrac{k}{x}$（$k > 0$），当 $x > 0$ 时，$y$ 随 $x$ 的增大而 \underline{\hspace{1cm}}；且当 $1 \leqslant x \leqslant 3$ 时，$y$ 的最大值与最小值的差为 $4$，则 $k$ 的值为（\quad）
\begin{enumerate}[label=\Alph*.]
\item $3$
\item $4$
\item $6$
\item $8$
\end{enumerate}

\end{enumerate}

\vspace{0.5em}

% ============================================================
% 二、填空题（共6题，每题3分，共18分）
% ============================================================
\noindent\textbf{二、填空题（本大题共6小题，每小题3分，共18分。）}

\begin{enumerate}[label=\arabic*., start=11]

\item 反比例函数 $y = \dfrac{8}{x}$ 的图象上有一点 $P(a, 2)$，则 $a =$ \underline{\hspace{3cm}}。

\item 若两个相似三角形的对应边之比为 $3 : 5$，则它们的面积之比为 \underline{\hspace{3cm}}。

\item 在 $\text{Rt}\triangle ABC$ 中，$\angle C = 90^\circ$，$AB = 10$，$AC = 6$，则 $\cos B =$ \underline{\hspace{3cm}}。

\item 一个正方体的棱长为 $3\,\text{cm}$，则该正方体的表面积为 \underline{\hspace{3cm}} $\text{cm}^2$。

\item 反比例函数 $y = \dfrac{k-1}{x}$ 的图象在第二、四象限，则 $k$ 的取值范围是 \underline{\hspace{3cm}}。

\item 如图，$AB \parallel CD$，$AC$ 与 $BD$ 交于点 $O$，$AO = 4$，$CO = 6$，$BO = 3$，则 $DO =$ \underline{\hspace{3cm}}。

\end{enumerate}

\vspace{0.5em}

% ============================================================
% 三、解答题（共8题，共72分）
% ============================================================
\noindent\textbf{三、解答题（本大题共8小题，共72分。解答应写出文字说明、演算步骤或证明过程。）}

\begin{enumerate}[label=\arabic*., start=17]

\item （8分）已知反比例函数 $y = \dfrac{k}{x}$ 的图象经过点 $A(2, 4)$。
\begin{enumerate}[label=(\arabic*)]
\item 求 $k$ 的值；
\item 判断点 $B(-4, -2)$ 是否在该函数图象上；
\item 当 $-4 \leqslant x \leqslant -1$ 时，求 $y$ 的取值范围。
\end{enumerate}
\vspace{6cm}

\item （8分）计算：
\begin{enumerate}[label=(\arabic*)]
\item $\sin^2 30^\circ + \cos^2 45^\circ + \tan 60^\circ$；
\item 在 $\text{Rt}\triangle ABC$ 中，$\angle C = 90^\circ$，$\angle A = 60^\circ$，$c = 8$，求 $a$ 和 $b$ 的值。
\end{enumerate}
\vspace{6cm}

\item （8分）如图，$\triangle ABC$ 中，$DE \parallel BC$，$D$ 在 $AB$ 上，$E$ 在 $AC$ 上。已知 $AD = 3$，$DB = 2$，$BC = 10$。
\begin{enumerate}[label=(\arabic*)]
\item 求 $\dfrac{DE}{BC}$ 的值；
\item 求 $DE$ 的长；
\item 若 $\triangle ADE$ 的面积为 $5.4$，求 $\triangle ABC$ 的面积。
\end{enumerate}
\vspace{6cm}

\item （9分）如图，某建筑物前有一根旗杆。小华站在距离建筑物底部 $12\,\text{m}$ 的点 $A$ 处，测得建筑物顶部的仰角为 $45^\circ$，旗杆顶部的仰角为 $60^\circ$（小华的身高忽略不计）。
\begin{enumerate}[label=(\arabic*)]
\item 求建筑物的高度；
\item 求旗杆的高度（结果保留根号）。
\end{enumerate}
\vspace{6cm}

\item （9分）已知反比例函数 $y_1 = \dfrac{k_1}{x}$ 和一次函数 $y_2 = k_2 x + b$ 的图象都经过点 $A(1, 6)$ 和点 $B(-2, m)$。
\begin{enumerate}[label=(\arabic*)]
\item 求反比例函数和一次函数的表达式；
\item 求 $\triangle AOB$ 的面积（$O$ 为坐标原点）；
\item 根据图象直接写出 $y_1 > y_2$ 时 $x$ 的取值范围。
\end{enumerate}
\vspace{6cm}

\item （10分）如图，某小区内有一块三角形空地 $ABC$。现要在空地内建一个矩形花园 $DEFG$，其中 $D$、$E$ 在 $BC$ 边上，$G$ 在 $AB$ 边上，$F$ 在 $AC$ 边上。已知 $BC = 24\,\text{m}$，$BC$ 边上的高 $AH = 8\,\text{m}$。
\begin{enumerate}[label=(\arabic*)]
\item 设矩形的宽 $DE = x\,\text{m}$，矩形的面积为 $S\,\text{m}^2$，写出 $S$ 与 $x$ 之间的函数关系式；
\item 当 $x$ 为多少时，矩形花园的面积最大？最大面积是多少？
\end{enumerate}
\vspace{6cm}

\item （10分）如图，某兴趣小组要测量一棵树 $CD$ 的高度。在距离树底部 $B$ 点 $20\,\text{m}$ 的 $A$ 点处，测得树顶 $D$ 的仰角为 $30^\circ$。然后向树的方向前进 $8\,\text{m}$ 到达 $E$ 点，再测得树顶 $D$ 的仰角为 $60^\circ$（测量者的身高忽略不计）。
\begin{enumerate}[label=(\arabic*)]
\item 求树的高度（结果保留根号）；
\item 若测量者的眼睛距地面 $1.5\,\text{m}$，则树的实际高度约为多少米？（参考数据：$\sqrt{3} \approx 1.73$，结果精确到 $0.1\,\text{m}$）
\end{enumerate}
\vspace{6cm}

\item （10分）如图，反比例函数 $y = \dfrac{k}{x}$（$x > 0$）的图象经过 $\text{Rt}\triangle OAB$ 的顶点 $A$，$AB \perp x$ 轴于点 $B$。已知 $OB = 4$，$\triangle OAB$ 的面积为 $6$。
\begin{enumerate}[label=(\arabic*)]
\item 求反比例函数的表达式；
\item 设点 $C$ 在反比例函数图象上，且 $\triangle OBC$ 的面积为 $3$，求点 $C$ 的坐标；
\item 在（2）的条件下，判断 $\triangle OAC$ 的形状，并求 $\angle AOC$ 的正切值。
\end{enumerate}
\vspace{6cm}

\end{enumerate}

\end{document}
